%! Inference and Experiments — Unit 3, Lesson 3.5 — Student Packet Cover Sheet
\documentclass[12pt]{article}
\usepackage{apstats-article}
\usepackage{apstats-boxes}
\usepackage{ltablex}
\keepXColumns

\begin{document}

% Full-bleed navy banner. \coverbanner MEASURES the title block and sizes the
% band to it, so a lesson title long enough to wrap grows the band rather than
% running out the bottom of it.
\coverbanner{Unit 3}{Lesson 3.5 \quad Inference and Experiments}

% NAMESTRIP: this is the ONLY component that carries the name row.
\namedateperiod
\vspace{0.18in}

\boxguard[14]
\begin{learningtargetbox}
\small
% GRADUAL RELEASE: the vocabulary is taught in the guided notes, so the targets name it
% plainly. The AP Learning Objective (VAR-3.E) stays in the teacher-facing plan.
\begin{enumerate}[leftmargin=1.5em, itemsep=5pt]
  \item \ldots{} read a simulation of random reshuffles and decide whether a difference
        between treatment groups is \textbf{statistically significant}.
  \item \ldots{} use the \textbf{scope of inference} to say whether a study shows cause and
        effect, and to whom its result can be \textbf{generalized}.
  \item \ldots{} explain why a random sample, without random assignment, cannot show cause
        and effect.
\end{enumerate}
\end{learningtargetbox}

\vspace{0.14in}

\boxguard[14]
\begin{tocbox}
\small
% A CONTENTS LIST ONLY: no Score column and no Total row. Whether a component counts is said in
% its description: AP Practice is extra credit, scored 1 to 5; the Homework is scored.
\renewcommand{\arraystretch}{1.5}
\begin{tabularx}{\linewidth}{@{} c l X @{}}
\rowcolor{sky}
\textbf{\#} & \textbf{Component} & \textbf{Description} \\
\midrule
1 & Warm-Up & Random selection or random assignment, breakfast and grades, and a music
                 experiment: could chance alone do it? \\
2 & Guided Notes \& Practice & Studying with music: a simulation, the scope of inference, and
                 cause or only association --- then a reading app together \\
3 & AP Practice & A sports drink, tomato plants, and green tea: AP-style multiple choice and
                 one free-response set --- extra credit, scored 1 to 5 \\
4 & Homework & A fruit-first lunch line and a city walking survey: significant or not, cause
                 or not, and to whom --- \textbf{scored 1 to 5, due the first class after two study halls} \\
\end{tabularx}
\end{tocbox}

\vspace{0.14in}

\boxguard[8]
\begin{remindbox}
\small
Even when a treatment does nothing, the luck of random assignment makes groups differ a
little. A difference is \textbf{statistically significant} only when chance alone would
rarely produce one that large. Then ask two separate questions: was the treatment
\emph{assigned} at random? That is what lets you say it \emph{caused} the difference. Were the
individuals \emph{selected} at random? That is what lets you speak for the whole population.
\textbf{A random sample is not a random assignment}: it lets you generalize, never conclude
cause.
\end{remindbox}

\end{document}
